\section{Base change and comparison of horizontal models}
\label{sec:comparison-v167}
We now distinguish the Hilbert graph of embedded curves from the
horizontal modification of a projective parameter family. In the
notation of Section~\ref{sec:horizontal-v166}, let $W\to B$ be
projective and of finite presentation, with $B$ integral of finite
type over $\C$. Fix a dense open $U\subset B$ over which $W$ is
flat with polynomial $P$. The admissible tests are locally
Noetherian morphisms $Y\to B$ for which $U_Y$ is schematically
dense. The horizontal object $W_Y^h$ is the schematic closure of
$W_{U_Y}$ in $W_Y$. Only tests for which this closure is flat
belong to the represented horizontal-flat subfunctor. Throughout,
\[
 H_B=\overline{s(U)}\ \subset\operatorname{Hilb}^{P}(W/B)
\]
means the \emph{horizontal Hilbert-main-component modification}.
Its terminal property refers to this admissible category and
this prescribed generic embedded quotient.

\begin{lemma}[Flat schematic density]\label{lem:density-v167}
Let $V\subset Y$ be a quasi-compact schematically dense open in a
locally Noetherian scheme. If $Z\to Y$ is flat, then $Z_V$ is
schematically dense in $Z$.
\end{lemma}
\begin{proof}
The assertion is local on $Y$ and $Z$. For $Y=\Spec A$, write
$V=\bigcup_{i=1}^rD(f_i)$. Schematic density is the injectivity
of $A\to\prod_i A_{f_i}$. Tensoring with any flat $A$-module
preserves this injection; finiteness of the product identifies
the tensor product with the product of the localizations. Apply
this to the rings of affine opens of $Z$. A section vanishing on
$Z_V$ is zero. This is precisely the required schematic density.
\end{proof}

\begin{theorem}[Strict base change and coherent comparison]
\label{thm:basechange-v167}
Let $B'\to B$ be an integral finite-type base change such that
$U'=U\times_BB'$ is nonempty. Use the pullback polarization and
write $H_{B'}$ for the horizontal Hilbert-main-component
modification of $W_{B'}$ relative to $U'$. There is a canonical
closed immersion
\begin{equation}\label{eq:comparison-closure-v167}
 H_{B'}=\overline{U'}^{\,H_B\times_BB'}
                    \lhook\joinrel\longrightarrow H_B\times_BB'.
\end{equation}
It is an equality if and only if $U'$ is schematically dense in
the full fibre product. In particular it is an equality for flat
base change. For composable base changes meeting the indicated
generic opens, these comparison maps compose to the comparison
for the composite. They obey the corresponding cocycle identity.

If $B^\nu\to B$ is the normalization, then $H_{B^\nu}\to H_B$
is finite and birational, and the normalizations of these two
integral schemes are canonically isomorphic over $B$.
\end{theorem}
\begin{proof}
The relative Hilbert scheme and its universal family commute with
base change. In the resulting Hilbert scheme over $B'$, the
closure of $s(U')$ is contained in the closed subscheme
$H_B\times_BB'$, since the ideal of $H_B$ vanishes on $s(U')$.
It is therefore exactly the schematic closure on the right of
\eqref{eq:comparison-closure-v167}. This proves the first
assertion and the density criterion. Flat base change preserves
the schematic density of $U\subset H_B$ by
Lemma~\ref{lem:density-v167}; note that the projection of the
fibre product to $H_B$ is flat.

For a further base change $B''\to B'$, both routes give the
schematic closure of $U''$ in the same relative Hilbert scheme.
The closed immersions consequently agree, including their scheme
structures. This proves transitivity and the cocycle, not merely
agreement on closed points. Finally, normalization of a finite-type
integral $\C$-scheme is finite. The map
$H_B\times_B B^\nu\to H_B$ is finite; its closed integral main
component is finite and birational because it agrees on a dense
normal open. Finite birational integral extensions have the same
integral closure in their common function field. This proves the
last assertion.
\end{proof}

\begin{example}[The full pullback need not commute]
\label{ex:basechange-v167}
Take $B=\A^2_{x,y}$, $W=\operatorname{Bl}_0B$, and
$U=B\setminus\{0\}$. The generic polynomial is $P=1$.
A flat family of length-one embedded subschemes of $W/B$ is a
section, so its relative Hilbert scheme is $W$ itself and
$H_B=W$. Restrict to $B'=\{y=0\}$. With homogeneous coordinates
$[u:v]=[x:y]$, the full pullback has equation $xv=0$ in
$\A^1_x\times\Pj^1$. It contains a vertical projective line.
The closure of $x\ne0$ is instead $v=0$, isomorphic to $B'$.
Thus \eqref{eq:comparison-closure-v167} is a strict inclusion.
The parameter $y=0$ is a nonflat base change; replacing the full
pullback by its reduction would not remove its extra component.
\end{example}

\begin{proposition}[Common refinements on one generic open]
\label{prop:refinement-v167}
Fix $B,U$. For two integral proper $U$-modifications $Y_1,Y_2$,
let $Y_1\star_UY_2$ be the schematic closure of the diagonal
$U$ in $Y_1\times_BY_2$. The projections are proper birational
maps and this object is terminal among integral $U$-modifications
equipped with maps to both $Y_1$ and $Y_2$ restricting to the
identity on $U$. Repeated products are canonically associative
and symmetric, and their comparison maps satisfy all product
cocycles. If a refinement of a stratum has the same reduced
closure and a smaller dense generic-flat open, its horizontal
model is canonically unchanged.
\end{proposition}
\begin{proof}
A pair of compatible maps gives a map into the product. Since its
integral source has dense $U$, the map factors through the
schematic closure of that open. The factorization is unique
because the closure is a closed subscheme. Properness follows
from properness of the product, and birationality from the common
dense open. For three or more factors every iterated object is
the closure of the same diagonal in the full product; this proves
associativity, symmetry, and their compatibility. Shrinking a
dense open of an integral section does not change its schematic
closure, proving the last assertion.
\end{proof}

The theorem supplies comparison for dominant restrictions and
strict base changes, and the proposition supplies common
refinements on a fixed generic problem. A stratum whose closure
lies entirely outside $U$ does not meet these hypotheses. In
particular these results do not identify Hilbert components with
different generic polynomials. For such strata the original
Noetherian-induction construction remains a choice of strata,
not a canonical incidence poset. If one asks for a flat family
whose fibres are the entire original fibres on two strata in one
connected base, equality of their Hilbert polynomials is necessary.
If $W$ is already flat over all of $B$, the horizontal model is
$B$ itself. These observations specify both the full-base flat
case and the precise limitation of the comparison theorem.

\subsection{Comparison with full flattening}
For an admissible $Y\to B$ on which $W_Y$ is flat,
Lemma~\ref{lem:density-v167} gives $W_Y=W_Y^h$. There is therefore
a natural inclusion of functors and a commutative diagram
\[
\begin{CD}
 \{Y:W_Y\text{ is flat}\}@>>>\{Y:W_Y^h\text{ is flat}\}\\
 @VVV @VVV\\
 \operatorname{Hom}_B(Y,B)@=\operatorname{Hom}_B(Y,B).
\end{CD}
\]
The top arrow is not in general surjective: the extra projective
line in Example~\ref{ex:basechange-v167} is exactly an obstruction
to full flattening. The usual flattening functor of a finitely
presented sheaf concerns its full tensor pullback and, when
represented, is represented by a monomorphism; see
\cite{Stacks166}, Tag~052F. By contrast, a horizontal modification
may have a positive-dimensional exceptional fibre. Flattening by
blow-up uses a pure transform and must likewise not be confused
with flatness of the full pullback. Theorem~\ref{thm:determinantal-v167}
computes a Rees ideal for the \emph{embedded-curve graph}; it does
not assert that the same ideal flattens every fibre of that
parameter graph.

The horizontal main component is integral, hence reduced. Neither
normality nor Cohen--Macaulayness follows from the construction:
for $W=B$ the construction returns $B$, which may be a nonnormal
integral curve or a non-Cohen--Macaulay integral higher-dimensional
scheme. The explicit surface in Theorem~\ref{thm:two-wall-v167}
has stronger properties because its charts have been computed.
Independence of polarization identifies representing objects for
the same horizontal functor; it does not identify the ambient
Hilbert schemes or their other components.
